
\documentclass{article}
\usepackage{colortbl}
\usepackage{makecell}
\usepackage{multirow}
\usepackage{supertabular}

\begin{document}

\newcounter{utterance}

\centering \large Interaction Transcript for game `eqbench', experiment `eq\_bench', episode 82 with reflexion{-}Qwen3.5{-}9B.
\vspace{24pt}

{ \footnotesize  \setcounter{utterance}{1}
\setlength{\tabcolsep}{0pt}
\begin{supertabular}{c@{$\;$}|p{.15\linewidth}@{}p{.15\linewidth}p{.15\linewidth}p{.15\linewidth}p{.15\linewidth}p{.15\linewidth}}
    \# & \multicolumn{2}{c}{Player} && \multicolumn{2}{c}{Game Master} \\
    \hline

    \theutterance \stepcounter{utterance}  
    & & & \multicolumn{4}{p{0.6\linewidth}}{
        \cellcolor[rgb]{0.9,0.9,0.9}{
            \makecell[{{p{\linewidth}}}]{
                \texttt{\tiny{[P1$\langle$GM]}}
                \texttt{You have attempted this same task before and did not succeed. These are the plans you wrote after those attempts.} \\
\texttt{Trial \#1: New plan: I failed because I assigned a high score to "Amused" (2) and "Humiliated" (3) without sufficient textual evidence. The dialogue is a tense, aggressive exchange where Charlie explicitly attacks Mari's competence ("too washed up") and asserts his seniority ("been on this beat longer"). There is no humor or self{-}deprecation in Charlie's lines to justify amusement or humiliation. My strategy should have been to strictly analyze the tone of Charlie's specific utterances: his defensive retort ("I've been on this beat longer") indicates anger and defensiveness, which maps directly to "Indignant." I should have assigned a score of 10 to Indignant and 0 to the others, as the context clearly depicts a power struggle rather than a comedic or shameful situation.} \\
\\ 
\texttt{Trial \#2: New plan: I failed because I assigned non{-}zero scores to "Amused" and "Humiliated" despite the dialogue containing no elements of humor or self{-}deprecation for Charlie. The exchange is a direct, aggressive power struggle where Charlie defends his seniority against Mari's attacks. My strategy must strictly map Charlie's defensive and aggressive tone ("I've been on this beat longer") to the "Indignant" label while assigning zero to all other emotions. I will output a score of 10 for Indignant and 0 for Amused, Humiliated, and Confused, ensuring no extraneous text is included.} \\
\\ 
\\ 
\\ 
\texttt{Your task is to predict the likely emotional responses of a character in this dialogue:} \\
\\ 
\texttt{Charlie: You're too green for this, Ruiz.} \\
\texttt{Mari: And you're too washed up, Redburn.} \\
\texttt{Charlie: I've been on this beat longer than you've been alive, kid.} \\
\texttt{Mari: And yet, here we are, stuck on the same case. Seems like experience isn't everything.} \\
\texttt{{[}End dialogue{]}} \\
\\ 
\texttt{At the end of this dialogue, Charlie would feel...} \\
\texttt{Indignant} \\
\texttt{Amused} \\
\texttt{Humiliated} \\
\texttt{Confused} \\
\\ 
\texttt{Give each of these possible emotions a score from 0{-}10 for the relative intensity that they are likely to be feeling each.} \\
\\ 
\texttt{You must output in the following format, including headings (of course, you should give your own scores), with no additional commentary:} \\
\\ 
\texttt{Indignant: <score>} \\
\texttt{Amused: <score>} \\
\texttt{Humiliated: <score>} \\
\texttt{Confused: <score>} \\
\\ 
\\ 
\texttt{{[}End of answer{]}} \\
\\ 
\texttt{Remember: zero is a valid score, meaning they are likely not feeling that emotion. You must score at least one emotion > 0.} \\
\\ 
\texttt{Your answer:} \\
            }
        }
    }
     \\ \\

    \theutterance \stepcounter{utterance}  
    & \multicolumn{4}{p{0.6\linewidth}}{
        \cellcolor[rgb]{0.561,0.737,0.561}{
            \makecell[{{p{\linewidth}}}]{
                \texttt{\tiny{[P1$\rangle$GM]}}
                \texttt{Indignant: 10} \\
\texttt{Amused: 0} \\
\texttt{Humiliated: 0} \\
\texttt{Confused: 0} \\
            }
        }
    }
    & & \\ \\

    \theutterance \stepcounter{utterance}  
    & & & \multicolumn{2}{p{0.3\linewidth}}{
        \cellcolor[rgb]{0.9,0.9,0.9}{
            \makecell[{{p{\linewidth}}}]{
                \texttt{\tiny{[GM$|$GM]}}
                \texttt{Indignant: 10} \\
\texttt{Amused: 0} \\
\texttt{Humiliated: 0} \\
\texttt{Confused: 0} \\
            }
        }
    }
    & & \\ \\

    \theutterance \stepcounter{utterance}  
    & & & \multicolumn{2}{p{0.3\linewidth}}{
        \cellcolor[rgb]{0.9,0.9,0.9}{
            \makecell[{{p{\linewidth}}}]{
                \texttt{\tiny{[GM$|$GM]}}
                \texttt{Indignant: 8} \\
\texttt{Amused: 0} \\
\texttt{Humiliated: 5} \\
\texttt{Confused: 0} \\
            }
        }
    }
    & & \\ \\

    \theutterance \stepcounter{utterance}  
    & & & \multicolumn{2}{p{0.3\linewidth}}{
        \cellcolor[rgb]{0.9,0.9,0.9}{
            \makecell[{{p{\linewidth}}}]{
                \texttt{\tiny{[GM$|$GM]}}
                \texttt{game\_result = LOSE} \\
            }
        }
    }
    & & \\ \\

    \theutterance \stepcounter{utterance}  
    & & \multicolumn{4}{p{0.6\linewidth}}{
        \cellcolor[rgb]{0.95,0.95,0.95}{
            \makecell[{{p{\linewidth}}}]{
                \texttt{\tiny{[P1$\leftrightarrow$P]}}
                \texttt{Post{-}episode Reflexion (not sent to GM):} \\
\texttt{New plan: I failed because I assigned non{-}zero scores to "Amused" and "Humiliated" despite the dialogue containing no elements of humor or self{-}deprecation for Charlie. The exchange is a direct, aggressive power struggle where Charlie defends his seniority against Mari's attacks. My strategy must strictly map Charlie's defensive and aggressive tone ("I've been on this beat longer") to the "Indignant" label while assigning zero to all other emotions. I will output a score of 10 for Indignant and 0 for Amused, Humiliated, and Confused, ensuring no extraneous text is included.} \\
\\ 
\texttt{Indignant: 10} \\
\texttt{Amused: 0} \\
\texttt{Humiliated: 0} \\
\texttt{Confused: 0} \\
            }
        }
    }
    & \\ \\

\end{supertabular}
}

\end{document}
